% ---------------------------------------------------------------------------
% SawitGuard-GNN --- Bab 3 "Solusi Usulan"
% Tempel ke badan paper. Butuh: \usepackage{graphicx}
% Figur: fig_pipeline.png (dihasilkan oleh ../fig_pipeline.py)
%
% Konvensi yang ditiru dari paper acuan: Bab 3 memuat RANCANGAN saja --- tidak ada
% satu pun angka HASIL MODEL. Angka yang muncul di sini seluruhnya adalah
% DESKRIPSI DATA (jumlah, jarak, derajat), bukan metrik performa; metrik masuk Bab 4.
%
% REVISI 2026-07-23: disusun ulang dari "dua lapisan terputus" menjadi SATU pipeline
% berkelanjutan enam tahap. Jembatan yang dulu digambar putus-putus kini terisi oleh
% Tahap 4 (rekonstruksi geometri kebun dari nama berkas ubin) --- lihat
% ../layer1_build/stitch_probe.py.
% ---------------------------------------------------------------------------

\section{Solusi Usulan}

Solusi yang dikembangkan adalah sebuah konsep sistem estimasi risiko penularan
antar-pohon pada perkebunan kelapa sawit, yang mengalirkan citra UAV mentah menjadi
angka peluang per-pohon: bagi setiap pohon yang \textit{hari ini masih tampak sehat},
berapa besar kemungkinan pohon tersebut tertular dalam beberapa siklus survei ke depan.
Pertanyaan yang dijawab bukan ``pohon mana yang sakit'' --- itu pekerjaan pengamatan,
bukan peramalan --- melainkan \textbf{ke mana penyakit akan bergerak selanjutnya, dan
apa sesungguhnya yang menentukan arah gerak itu}. Rumusan tersebut dipilih karena hanya
inilah pertanyaan yang dapat dijawab secara sahih dengan data yang tersedia, dan karena
keluarannya, yaitu daftar pohon berperingkat risiko, adalah bentuk yang langsung dapat
dipakai tim kebun untuk mengarahkan inspeksi darat alih-alih menyisir seluruh blok.

Pipeline disusun sebagai enam tahap berurutan, dari piksel sampai peluang, dan dirangkum
pada Gambar~\ref{fig:pipeline}. Tiga tahap pertama bekerja pada citra UAV nyata dan
menghasilkan inventaris tajuk. Tahap keempat merekonstruksi tata letak kebun yang
sebenarnya dari inventaris tersebut. Dua tahap terakhir menjalankan dinamika penularan
di atas tata letak nyata itu dan menerjemahkan keluaran model menjadi persentase risiko
yang terkalibrasi.

\begin{figure}[h]
  \centering
  \includegraphics[width=\linewidth]{fig_pipeline.png}
  \caption{Pipeline enam tahap SawitGuard-GNN, dari citra UAV sampai persentase risiko
  per-pohon. Tahap 1--4 berjalan pada data nyata; Tahap 5 menjalankan dinamika penularan
  sintetis di atas geometri kebun yang nyata; Tahap 6 mengubah skor model menjadi peluang
  terkalibrasi. Batas antara yang nyata dan yang disimulasikan ditandai eksplisit dan
  dirinci pada Tabel~\ref{tab:data}.}
  \label{fig:pipeline}
\end{figure}

\subsection{Tahap 1--3 --- Dari Citra UAV Menjadi Inventaris Tajuk}

Tiga tahap pertama mengubah citra UAV RGB bersudut nadir menjadi daftar tajuk sawit
yang masing-masing disertai posisi, ukuran, dan skor kesehatan.

\textbf{Deteksi tajuk} memakai YOLO11 untuk menemukan setiap tajuk sawit di dalam ubin
citra dan memberinya kotak pembatas. Keluaran tahap ini adalah koordinat piksel pusat
tajuk beserta lebar dan tingginya --- dan koordinat inilah, bukan sekadar hitungan
jumlah pohon, yang kelak menjadi tulang punggung Tahap~4.

\textbf{Estimasi luas tajuk} mengambil setiap kotak hasil deteksi dan memisahkan piksel
tajuk dari piksel tanah memakai indeks \textit{Excess Green} yang diambangkan dengan
metode Otsu. Pendekatan ini sengaja dipilih tanpa pelatihan, sebab data yang tersedia
hanya menyediakan kotak pembatas dan \textbf{tidak memiliki \textit{mask}
kebenaran-dasar}. Konsekuensinya dinyatakan di muka dan tidak dilunakkan: tahap ini
diperlakukan sebagai demonstrasi kualitatif belaka, tidak ada klaim IoU yang dapat
dipertanggungjawabkan, dan karena itu tidak ada yang diajukan. Luas tajuk dipakai
sebagai proksi ukuran kanopi, yang berguna sebagai penciri ukuran pohon pada tahap
berikutnya, bukan sebagai hasil segmentasi yang divalidasi.

\textbf{Penilaian kesehatan tajuk} menghitung statistik warna dan tekstur kanal RGB di
dalam setiap kotak --- simpangan baku kanal merah, selisih hijau--merah, rerata kanal
biru, selisih hijau--biru --- lalu mengklasifikasikannya dengan LightGBM. Pemilihan
fitur ini disengaja dan bukan hasil pencarian buta: klorosis, yaitu gejala visual yang
membuat tajuk menguning, memang termanifestasi sebagai pergeseran kehijauan, sehingga
fitur yang dipilih memiliki alasan agronomis sebelum diuji secara empiris.

\subsection{Tahap 4 --- Rekonstruksi Geometri Kebun}

Tahap keempat adalah penghubung yang membuat pipeline ini berkelanjutan, dan merupakan
tambahan utama pada rancangan ini. Persoalannya sederhana namun menentukan: model
penularan membutuhkan \textbf{peta posisi pohon}, sedangkan Tahap~1 hanya menghasilkan
koordinat di dalam ubin citra selebar 1.024 piksel yang saling terpisah.

Kunci pemecahannya terletak pada penamaan berkas ubin, yang menyimpan offset piksel
absolut ubin tersebut di dalam ortomosaik induknya. Dengan menjumlahkan offset ubin dan
posisi pusat kotak pembatas, setiap tajuk memperoleh \textbf{koordinat global} di dalam
ortomosaiknya. Seluruh ubin dengan demikian dapat dijahit kembali menjadi satu awan
titik yang merepresentasikan tata letak kebun yang sebenarnya.

Satu langkah pembersihan wajib menyertainya. Ubin-ubin pada data ini beririsan sangat
rapat, sehingga satu pohon yang sama muncul berulang kali pada ubin yang berbeda-beda.
Deteksi ganda tersebut digabungkan dengan aturan jarak: dua deteksi yang terpisah kurang
dari setengah ukuran tajuk median dianggap pohon yang sama. Besarnya penggabungan yang
terjadi --- dan implikasinya terhadap ukuran sampel efektif --- dilaporkan apa adanya
pada Subbab~3.6, sebab hal itu mengubah cara seluruh hasil Lapisan~1 harus dibaca.

Awan titik hasil rekonstruksi lalu dikonversi ke satuan meter. Faktor konversinya
diperoleh dengan mengukur jarak tetangga terdekat antar pohon dan memadankannya dengan
jarak tanam sawit yang lazim, yaitu sekitar 9~m. Konsistensi hasilnya menjadi
pemeriksaan-silang yang berdiri sendiri: ketiga ortomosaik, yang diproses secara
terpisah, menghasilkan skala yang saling berdekatan, sehingga asumsi jarak tanam
tersebut tidak bertentangan dengan datanya sendiri.

\subsection{Tahap 5 --- Graf Multi-relasi dan Dinamika Penularan}

Di atas awan titik nyata hasil Tahap~4 dibangun sebuah \textbf{graf multi-relasi} dengan
tiga relasi yang mewakili tiga jalur penularan yang berbeda sifatnya. Relasi
\textit{akar} menghubungkan pohon-pohon yang bersentuhan perakaran dalam radius pendek
dan merupakan jalur terkuat, sesuai pemahaman bahwa \textit{Ganoderma} terutama menyebar
lewat kontak akar. Relasi \textit{angin} berbentuk kerucut berarah mengikuti angin
dominan, menjangkau lebih jauh dengan tekanan jauh lebih lemah. Relasi
\textit{manajemen} menghubungkan pohon-pohon di sepanjang jalur panen yang sama, mewakili
penularan yang terbawa alat dan lalu lintas pekerja.

Karena graf kini dibangun di atas posisi pohon yang sebenarnya, kerapatan sambungannya
bukan lagi angka yang dipilih perancang melainkan sifat yang \textbf{terukur dari data}.
Nilainya dilaporkan pada Tabel~\ref{tab:data} dan berbeda dari kisi ideal --- kebun nyata
memiliki jarak tanam yang tidak seragam, jalan, dan rumpang tanaman, sehingga rata-rata
jumlah tetangga sebuah pohon lebih rendah daripada kisi teratur. Perbedaan ini tidak
diperhalus; sebaliknya, kepekaan hasil terhadap kerapatan sambungan diuji secara
eksplisit sebagai salah satu ablasi pada Bab~4.

Dinamika penularannya sendiri dibangkitkan dengan model SEIR berjaring
(\textit{Susceptible}, \textit{Exposed}, \textit{Infected}, \textit{Removed}) selama 20
siklus survei, dengan satu siklus mewakili kira-kira tiga bulan. Di sinilah batas
kejujuran pipeline ini harus ditarik dengan tegas, dan ia ditarik di muka: \textbf{tata
letak kebunnya nyata, tetapi dinamika penularannya disimulasikan}. Seluruh laju epidemi
--- laju penularan tiap relasi, laju inkubasi $E \rightarrow I$, dan laju $I \rightarrow R$
--- dipusatkan pada satu berkas konfigurasi dan ditandai eksplisit sebagai
\textit{placeholder} yang wajib dikalibrasi ke literatur \textit{Ganoderma} sebelum klaim
lapangan apa pun diajukan. Alasannya bukan kelalaian melainkan ketiadaan data: peramalan
penyebaran menuntut deret waktu per-pohon dengan status infeksi berlabel, dan data
semacam itu tidak tersedia untuk BSR di ranah publik.

Fitur node sengaja tidak dibangkitkan secara acak. Setiap pohon memancarkan spektrum
sesuai keadaannya, diambil dari \textbf{data hiperspektral penyakit layu pinus yang
nyata}: keadaan $S$ dan $I$ diambil langsung dari data tersebut, keadaan $E$ dibentuk
sebagai interpolasi $E = (1-\alpha)S + \alpha I$, dan keadaan $R$ diambil dari spektrum
terinfeksi paling ekstrem. Nilai $\alpha$ dikunci pada 0{,}1, yaitu satu-satunya nilai
yang membuat keterpisahan $S$ terhadap $E$ berada pada rentang yang menantang namun tidak
mustahil --- apabila $E$ terlalu mudah dibedakan dari $S$, tugas peringatan dini menjadi
sepele dan kehilangan makna. Spektrum lalu direduksi dengan PCA dan disatukan dengan
selisih temporalnya menjadi vektor fitur per pohon per siklus. Dua pengaman melekat di
sini: penyekala dan PCA \textbf{hanya dipasang pada siklus latih}, dan selisih temporal
hanya menengok ke siklus sebelumnya, sehingga tidak ada informasi masa depan yang bocor
ke masukan model. Batasnya tetap dinyatakan terbuka --- spektrum ini berasal dari penyakit
layu pinus, bukan BSR, dan kompartemen $E$ merupakan interpolasi yang tidak ada dalam data.

\subsection{Tahap 6 --- Model Peramalan dan Estimasi Risiko Terkalibrasi}

Tugas peramalan dirumuskan sebagai pertanyaan peringatan dini: bagi sebuah pohon yang
masih \textbf{tanpa gejala} pada siklus $t$, apakah pohon tersebut akan menjadi bergejala
dalam $h$ siklus ke depan? Pohon yang sudah bergejala pada siklus $t$ dikeluarkan dari
kumpulan contoh, karena meramalkannya bukan peringatan dini melainkan pencatatan. Setiap
model menerima jendela tiga siklus terakhir dan diuji pada tiga cakrawala waktu.
Pembagian data bersifat temporal, siklus awal untuk melatih dan siklus akhir untuk menguji.

Tiga model diadu dengan tugas yang persis sama. \textbf{MLP} bekerja per-pohon tanpa graf
sama sekali dan berperan sebagai garis dasar lapangan. \textbf{STGNN} menambahkan difusi
multi-relasi dengan bobot perhatian yang dilatih untuk ketiga relasi, lalu meringkas
riwayatnya dengan GRU. \textbf{STGNN-SEIR} identik dengan STGNN kecuali satu tambahan:
sebuah kepala SEIR yang dapat dilatih, yang membaca keadaan lunak $(S,E,I)$ dari embedding
graf dan menduga peluang bergejala memakai laju yang dipelajari serta suku tekanan
tetangga. Satu pengaman berlaku bagi ketiganya: \textbf{tidak ada model yang pernah
melihat keadaan SEIR sebenarnya sebagai masukan} --- yang masuk hanyalah fitur spektral
dan agregat tetangganya. Garis dasar MLP sengaja disetel dengan baik, sebab pertanyaannya
adalah apakah graf mengungguli model per-pohon yang benar-benar bagus, bukan model lemah
yang dipasang untuk dikalahkan.

Keluaran mentah ketiga model berupa skor yang hanya sahih sebagai \textbf{peringkat}.
Skor semacam itu cukup untuk mengurutkan pohon dari paling berisiko ke paling aman, namun
belum boleh dibaca sebagai persentase --- menyebut sebuah pohon ``berisiko 23\%'' menuntut
bahwa dari seratus pohon yang diberi angka serupa, kira-kira dua puluh tiga di antaranya
benar-benar tertular. Tahap terakhir karena itu \textbf{mengkalibrasi} skor menjadi
peluang, dengan pemeta kalibrasi yang \textbf{hanya dipasang pada siklus latih} dan
diterapkan apa adanya pada siklus uji. Mutu kalibrasi dilaporkan sebagai kurva keandalan
beserta skor Brier, sehingga pembaca dapat menilai sendiri apakah angka persentase yang
dikeluarkan sistem memang berarti seperti yang dituliskannya.

Peluang terkalibrasi itu kemudian diringkas menjadi bentuk yang dapat dipakai di lapangan.
Pohon dikelompokkan ke dalam tingkat risiko, dan pada tiap tingkat dilaporkan presisi
pada-$k$-teratas: apabila tim kebun hanya sanggup memeriksa sejumlah pohon dalam satu
siklus, berapa banyak di antaranya yang memang akan tertular. Metrik inilah yang
menjembatani keluaran statistik dengan keputusan operasional, dan ia dilaporkan
berdampingan dengan \textit{lead time} --- berapa siklus lebih awal peringatan muncul
sebelum gejala tampak --- sebab sebuah model dapat unggul pada peringkat namun tidak
memberi keuntungan waktu apa pun di lapangan.

\subsection{Data}

Seluruh sumber data, jumlahnya, dan status nyata atau sintetisnya dirangkum pada
Tabel~\ref{tab:data}. Tabel ini sengaja mencantumkan kedua angka jumlah tajuk --- jumlah
anotasi mentah dan jumlah pohon unik setelah rekonstruksi Tahap~4 --- karena selisih
keduanya sangat besar dan menentukan bagaimana hasil Lapisan~1 harus dibaca.

\begin{table}[h]
\centering
\caption{Sumber data tiap tahap, beserta status nyata atau sintetisnya. Angka tajuk
dilaporkan dalam dua bentuk: anotasi mentah, dan pohon unik setelah deteksi ganda
antar-ubin digabungkan pada Tahap~4.}
\label{tab:data}
\begin{tabular}{llll}
\hline
\textbf{Tahap} & \textbf{Sumber dan jumlah} & \textbf{Label} & \textbf{Status} \\
\hline
1--3. Inventaris tajuk & 3 ortomosaik UAV RGB nadir; 2.303 ubin 1.024 px; 151.060 kotak anotasi & Healthy / Unhealthy (kesehatan tajuk generik, bukan BSR) & Nyata \\
4. Geometri kebun & 5.077 pohon unik (1.379 / 1.849 / 1.849 per ortomosaik); 96,6\% anotasi adalah deteksi ganda antar-ubin & Posisi global per pohon & Nyata (diturunkan) \\
4. Skala dan luas & Jarak tetangga terdekat median 101--106 px; skala 8,5--8,9 cm/px; luas blok 16,9 / 17,6 / 18,8 ha & Jarak tanam acuan 9 m (asumsi kalibrasi) & Nyata (terukur) \\
5. Struktur graf & Derajat relasi akar terukur 5,37--5,52 pada radius 13 m & Tiga relasi: akar, angin, manajemen & Nyata (topologi) \\
5. Dinamika penularan & SEIR berjaring, 20 siklus survei, 1 siklus $\approx$ 3 bulan & Status S/E/I/R per pohon per siklus & Sintetis (laju placeholder) \\
5. Fitur node & Spektra hiperspektral penyakit layu pinus & Spektrum per keadaan; $E$ interpolasi $\alpha = 0,1$ & Nyata tetapi penyakit lain \\
6. Kalibrasi risiko & Pemeta kalibrasi dipasang pada siklus latih saja & Peluang tertular dalam $h$ siklus & Diturunkan dari Tahap 5 \\
\hline
\end{tabular}
\end{table}

Pembacaan yang jujur atas tabel ini adalah sebagai berikut. Segala sesuatu yang menyangkut
\textbf{ruang} --- di mana pohon berdiri, seberapa rapat, siapa bertetangga dengan siapa,
berapa luas bloknya --- bersumber dari data nyata dan terukur. Segala sesuatu yang
menyangkut \textbf{waktu} --- kapan sebuah pohon tertular, seberapa cepat penyakit
menjalar, berapa lama masa inkubasinya --- disimulasikan. Pipeline ini karena itu dapat
menguji apakah sebuah model peramalan benar-benar memanfaatkan struktur ruang yang nyata,
tetapi belum dapat menguji apakah penyebaran BSR yang sesungguhnya mengikuti dinamika yang
diasumsikan.

\subsection{Protokol Evaluasi Anti-Kebocoran}

Protokol evaluasi ditetapkan sebelum hasil dilihat, dan berbeda untuk tiap tahap karena
sumber kebocorannya berbeda.

Pada \textbf{Tahap 1--3} hanya tersedia tiga ortomosaik. Rekonstruksi Tahap~4 mengungkap
bahwa ubin-ubin di dalam satu ortomosaik beririsan jauh lebih rapat daripada yang tampak:
dari 151.060 kotak anotasi hanya tersisa 5.077 pohon unik, sehingga rata-rata setiap pohon
tercatat sekitar tiga puluh kali pada ubin yang berlainan. Pembagian acak karena itu tidak
sekadar berisiko bocor, melainkan \textbf{pasti bocor}: ubin yang memuat pohon yang sama
akan tersebar ke sisi latih dan sisi uji sekaligus. Seluruh evaluasi karena itu memakai
\textbf{validasi silang blok \textit{leave-one-ortho-out}}, yaitu setiap lipatan menahan
satu ortomosaik utuh, dan asal ortomosaik dibaca dari nama berkas tiap ubin.

Pengulangan tiga puluh kali lipat itu membawa satu konsekuensi lagi yang harus dinyatakan
dan tidak boleh disembunyikan di balik jumlah baris yang besar: \textbf{daya statistik
sesungguhnya jauh lebih kecil daripada yang disiratkan angka 151.060}. Pada tingkat pohon
unik, kelas \textit{Unhealthy} hanya berjumlah \textbf{66 pohon} di seluruh tiga ortomosaik
--- 17, 31, dan 18 pohon pada masing-masing lipatan. Setiap selang kepercayaan karena itu
harus dihitung pada tingkat pohon unik, bukan tingkat anotasi, dan setiap hasil Lapisan~1
harus dibaca sebagai berbasis sampel kecil. Metriknya
adalah mAP untuk deteksi serta PR-AUC dan ROC-AUC untuk kesehatan. Akurasi tidak
dilaporkan: pada ketimpangan kelas seekstrem ini, pengklasifikasi yang selalu menjawab
\textit{Healthy} sudah mencapai akurasi di atas 98\% tanpa memiliki kegunaan apa pun.

Pada \textbf{Tahap 5--6} setiap konfigurasi dijalankan pada 20 \textit{seed} dan
dilaporkan sebagai rerata $\pm$ simpangan baku. Perbandingan antarmodel dilakukan
\textbf{berpasangan per-\textit{seed}}, disertai hitungan tanda yang menunjukkan pada
berapa dari 20 \textit{seed} sebuah model unggul --- sebab menang tipis pada 20 dari 20
\textit{seed} dan menang besar pada 11 dari 20 adalah dua pernyataan yang sangat berbeda.
Aturan keputusannya ditetapkan di muka: apabila nilai mutlak rerata selisih lebih kecil
daripada simpangan bakunya, hasil dinyatakan \textbf{tidak konklusif} dan tidak pernah
dibulatkan menjadi kemenangan.

\subsection{Rancangan Ablasi --- Apa yang Sebenarnya Menentukan Risiko}

Selisih ``STGNN dikurangi MLP'' saja tidak dapat ditafsirkan, karena mencampur tiga efek
yang berbeda menjadi satu angka. Rancangan ablasi karena itu memisahkannya dengan cara
menukar \textbf{tampilan graf} yang dilihat model sambil menahan arsitekturnya tetap sama:
tampilan \textit{asli} berisi graf kontak yang dibangun di atas geometri kebun nyata,
tampilan \textit{tanpa-graf} menihilkan seluruh sisi, tampilan \textit{acak}
mempertahankan derajat yang sama tetapi merusak strukturnya, dan tampilan
\textit{perturbasi} menginterpolasi keduanya menurut parameter $\varepsilon$. Ketiga
selisih berikut bersifat aditif dan menyusun keseluruhan selisih STGNN--MLP:

\begin{itemize}\itemsep2pt
  \item \textbf{temporal} $=$ tanpa-graf $-$ MLP --- nilai dari sekadar memodelkan
        riwayat pohon itu sendiri;
  \item \textbf{prevalensi} $=$ acak $-$ tanpa-graf --- nilai dari mengetahui
        \textit{seberapa banyak} penyakit ada di sekitar, tanpa mengetahui di mana;
  \item \textbf{struktur} $=$ asli $-$ acak --- nilai dari mengetahui \textit{siapa
        persis} tetangga sebuah pohon.
\end{itemize}

Dekomposisi inilah inti ilmiah paper ini, sebab ia menjawab langsung pertanyaan yang
diajukan pada awal bab: apa yang sesungguhnya menentukan risiko penularan sebuah pohon.
Komponen \textbf{struktur} adalah yang mempertaruhkan premis pendekatan berbasis graf.
Apabila peta kontak multi-relasi memang penting, komponen inilah yang harus besar;
apabila kecil, keunggulan model graf sesungguhnya berasal dari pemodelan waktu dan
tekanan wabah agregat, dan hal itu wajib dinyatakan sebagai temuan, bukan disembunyikan
di balik selisih agregat. Dekomposisi ini dirancang justru untuk membuka kemungkinan
tersebut, dan kesimpulannya dilaporkan pada Bab~4 apa adanya, ke arah mana pun ia jatuh.

Dua ablasi pendamping menguji ketahanan kesimpulan itu. Yang pertama merusak peta kontak
secara bertahap menurut $\varepsilon$ untuk melihat pada tingkat kesalahan peta berapa
keunggulan graf lenyap. Yang kedua menyapu kerapatan sambungan relasi akar, dan justru di
sinilah geometri nyata hasil Tahap~4 memperoleh perannya yang paling menentukan: derajat
yang terukur dari kebun sebenarnya harus berada di dalam rentang yang disapu, sebab hanya
dengan begitu kesimpulan yang ditarik dari simulasi dapat diklaim berlaku pada kerapatan
tanam yang nyata.
